% Section 8.7, from lectures/L08/README.md.

\section{Summary}
\label{c8:sec:review}

\needspace{6\baselineskip}
\subsection*{What you should be able to do now}
\begin{itemize}
  \item Explain what a finite state machine is, and point out its states, transitions, inputs and
    outputs in a given circuit.
  \item Design a Moore machine by hand, from specification through state diagram and state table to
    Karnaugh-derived equations, and build the result in CircuitVerse.
  \item Model the same machine in VHDL with an enumerated type \code{state_t} and a \code{case}
    process.
  \item Recognise a Mealy machine on sight and say what its output costs and buys, one clock cycle
    of latency against one extra state, and why Moore is the default here and in the CAN project.
  \item Compose a state machine with previously built subcomponents instead of rewriting them. In
    the capstone that means three earlier chapters' modules around a machine that is the only new
    logic, with shift and counter idioms from one more.
\end{itemize}

\needspace{6\baselineskip}
\subsection*{Questions to test yourself with}
\begin{enumerate}
  \item What makes a circuit a state machine, rather than just a register with some logic around
    it?
  \item What exactly is \code{X} in the hand-designed machine, and why does it have to be a
    one-clock-cycle pulse rather than the button's level?
  \item Why is that machine's state encoding a Gray code? What can go wrong with an ordinary binary
    count?
  \item A counter and a shift register are technically state machines too. What separates them from
    the machine designed here?
  \item Which signals may a Moore output depend on, and which may a Mealy output depend on? Given a
    single line of VHDL driving an output, how do you decide which of them you are looking at?
  \item Why does the Moore version in \secref{c8:sec:mealy}'s comparison need an extra state, and
    why is its output exactly one clock cycle behind the Mealy version's?
  \item Why does \code{LED_PROCESS} in \file{fsm_led.vhd} read only \code{state}, never
    \code{button_edge_s2} or \code{button_n} directly? What would change if it did?
  \item Why does \code{din} in \code{seq_detect_mealy} not go through a synchroniser, unlike
    \code{reset_n}?
  \item What does \code{when others} buy you in a \code{case} for state transitions, when
    \code{state_t} only has the values you declared?
  \item Why does the timer in the capstone receiver run sixteen times as fast as the baud rate
    rather than once per bit, and why does sampling mid-bit make the design tolerate a transmitter
    whose clock differs slightly from yours?
  \item The capstone transmitter uses a full-bit timer and needs nothing mid-bit. What does it have
    that the receiver does not, which makes the difference?
\end{enumerate}

\needspace{6\baselineskip}
\subsection*{Looking ahead}
This is the last chapter of the book's first part. Over eight chapters we have gone from gates and
Karnaugh maps, by way of registers, metastability and the VHDL language itself, to counters, shift
registers and timers, and now finite state machines, at every step reasoning about the circuit by
hand before writing it in synthesisable VHDL and running it on a real FPGA. Every pattern built
along the way, the two-flip-flop synchroniser, the \code{process} and \code{case} coding style, and
now the Moore machine, has been reused in this chapter's own examples, and carries over unchanged
into what comes next.

Next is \chapref{c9:ch}, the first practical exam: three hours, individual, at the computer with
GHDL, over everything the first eight chapters have built. After it the group project begins, where
a CAN controller is designed in VHDL over ten chapters and finishes against a real microcontroller
over SPI.
